\documentclass[11pt,a4paper]{article}

\usepackage[T1]{fontenc}
\usepackage[utf8]{inputenc}
\usepackage{lmodern}
\usepackage[english]{babel}
\usepackage{amsmath,amssymb,amsthm,mathtools}
\usepackage{geometry}
\usepackage{booktabs}
\usepackage{enumitem}
\usepackage{graphicx}
\usepackage{xcolor}
\usepackage{hyperref}
\usepackage{microtype}

\geometry{margin=2.4cm}
\hypersetup{colorlinks=true,linkcolor=black,citecolor=black,urlcolor=blue}

\theoremstyle{definition}
\newtheorem{problem}{Problem}
\newtheorem{definition}{Definition}

\newcommand{\todo}[1]{\textcolor{red}{[TODO: #1]}}

\providecommand{\keywords}[1]{\par\medskip\noindent\textbf{Keywords:} #1}

\title{Cooperative Spectrum Reservation for Datacenter Migration\\
in Multi-Provider Elastic Optical Networks\\
during the Disaster Warning Window}
\author{Cláudio de Faria Dantas}
\date{\todo{date}}

\begin{document}
\maketitle

\begin{abstract}
When a disaster is forecast, operators have a short warning window to move critical data away from the infrastructure that is about to fail.
In an elastic optical network shared by several providers, each provider sees only its own overlay, learns about the threat at its own time and allocates spectrum on its own, while all of them compete for the same contiguous slots on the links that leave the threatened site.
We study whether providers that agree to cooperate can migrate more data before the failure than providers that act alone.
Cooperation is restricted to exchanging information: no provider allocates spectrum for another.
We propose \emph{cooperative corridor protection}: on the links of a cooperator's migration path that, according to a per-link tenant count disclosed by the infrastructure owner, are used only by cooperating providers, background traffic is kept out of a reserved part of the spectrum until the failure.
The reservation is either a fixed fraction of the link or sized by the migration still pending.
We evaluate the mechanism with a discrete-event simulator of a 24-node backbone shared by five providers, over every possible cooperating subset, against disaster-aware First-Fit without reservation and against providers that protect only the links they use alone.
\todo{one or two sentences with the main quantitative results once the final runs are in.}
\end{abstract}

\keywords{elastic optical networks $\cdot$ disaster preparedness $\cdot$ datacenter migration $\cdot$ multi-provider cooperation $\cdot$ spectrum reservation}

%======================================================================
\section{Introduction}
%======================================================================

Earthquakes, floods, wildfires and large power outages damage communication infrastructure in geographically correlated ways: many nodes and links fail together, and the datacenters attached to them become unreachable \cite{rak2021disaster}.
Most work on disaster resilience in optical networks addresses what happens \emph{after} the failure, restoring as many connections as possible with the surviving resources \cite{zou2022drama,zou2025dramapro,sahoo2022cooperation}.
Some disasters, however, can be forecast minutes to hours in advance.
In that \emph{warning window}, the network is still intact and the most valuable action is preparedness: moving critical data out of the area that is about to fail.

This work focuses on that window in an elastic optical network (EON) whose physical infrastructure is shared by several providers.
Each provider operates a virtual overlay of the physical topology, sees only its own overlay, discovers the threat at its own time and allocates spectrum independently.
All of them hold critical data at the threatened site and must migrate it before the failure, through links that are partly shared and whose spectrum must be contiguous and identical along each path.
The links leaving the threatened site are typically shared by several providers, and migration competes there with the background traffic that keeps flowing during the warning window.

The question we ask is whether providers that \emph{cooperate} can complete more of their migration than providers that act alone, under the constraint that cooperation is limited to exchanging information.
There is no central allocator and no provider allocates spectrum for another.

\paragraph{Contributions.}
\begin{itemize}[leftmargin=1.2em]
  \item A formulation of the multi-provider migration problem in the disaster warning window, with independent action, independent discovery and information-only cooperation (Section~\ref{sec:problem}).
  \item \emph{Cooperative corridor protection}, a mechanism in which each cooperator restricts its own background traffic on the links of its migration path that the cooperating group can prove are used by no other provider. Two reservation sizes are studied: a fixed fraction of the link and a demand-driven size that follows the migration still pending (Section~\ref{sec:mechanism}).
  \item An evaluation over every subset of cooperating providers, comparing cooperative protection with disaster-aware First-Fit without reservation and with an individual variant in which each provider protects only the links it uses alone (Sections~\ref{sec:setup} and~\ref{sec:results}).
\end{itemize}

\todo{organization paragraph once the final section list is fixed.}

%======================================================================
\section{Background and related work}
\label{sec:background}
%======================================================================

\subsection{Elastic optical networks}

In an EON, the spectrum of each fiber is divided into narrow frequency slots, and a connection (lightpath) occupies a block of adjacent slots sized to its bit rate \cite{jinno2009spectrum}.
Routing and spectrum assignment (RSA) must respect two constraints: \emph{contiguity}, the slots of a connection are adjacent, and \emph{continuity}, the same slots are used on every link of the path.
The number of slots also depends on the modulation format, which is limited by the transmission reach: shorter paths can use more efficient formats and therefore fewer slots.
First-Fit, which tries candidate paths in order and takes the lowest-indexed block of free slots that fits, is a common baseline because of its simplicity and low fragmentation.

\subsection{Disasters and the phases of resilience}

Disasters differ from isolated failures because they affect many geographically correlated components at once, which invalidates the independence assumptions behind classical protection schemes such as dedicated or shared backup path protection \cite{shen2016survivable,rak2021disaster}.
A common abstraction models the disaster zone as a region in which all nodes and incident links fail simultaneously \cite{zou2022drama,zou2025dramapro}.
Resilience is usually organized in phases: preparedness before the event, response during it and recovery afterwards \cite{rak2021disaster}.
This work belongs to preparedness, more precisely to the interval between the warning and the failure.

\subsection{Post-disaster restoration in EONs}

A large body of work addresses restoration after the failure.
The DRAMA family of algorithms \cite{zou2022drama,zou2025dramapro} introduces a mitigation zone around the disaster area, in which surviving lightpaths can be degraded in a controlled way to free spectrum and increase the number of recovered demands, combined with penalty-aware routing and adaptive modulation.
These algorithms assume a single operator with a global view of the network and act after the failure; our problem has several providers with local views and takes place before it.

\subsection{Cooperation across administrative domains}

Disasters affect the transport network and the datacenters that depend on it at the same time, which has motivated cooperation between datacenter providers and optical carriers.
Sahoo et al.\ \cite{sahoo2022cooperation} show that a mediator through which datacenter providers and carriers share abstracted topology information and jointly decide how to use surviving resources substantially improves recovery compared with non-cooperative operation.
Our setting shares the idea that cooperation is mediated and limited in what it reveals, but differs in three ways: the actors are providers sharing one optical infrastructure, the action takes place in the warning window, and cooperation never transfers the right to allocate spectrum.

\subsection{Information disclosed by infrastructure owners}

Our mechanism assumes that the owner of the optical infrastructure discloses, for each link, how many providers use it.
Related disclosures already exist.
Carriers tag links with shared risk link groups (SRLGs) so that customers can request paths that do not share a fiber or duct, and their collection along a path is standardized \cite{rfc8001}.
In spectrum-as-a-service offerings, the owner of an open line system knows every tenant of every link.
A per-link tenant count, revealing a number but not identities, is a smaller disclosure than either; still, it is an assumption of this work and not an established practice.

%======================================================================
\section{Problem formulation}
\label{sec:problem}
%======================================================================

\subsection{Physical network and spectrum}

The physical topology is an undirected weighted graph $G=(V,E)$ with link length $\ell:E\to\mathbb{R}_{>0}$.
Each link $e$ carries $F$ frequency slots, with occupancy vector $\mathbf{s}_e(t)\in\{0,1\}^F$.

\begin{definition}[Spectral allocation]
An allocation for a request $r$ of bit rate $b_r$ is a pair $(p,[f_\mathrm{s},f_\mathrm{e}])$, where $p$ is a simple path and
\begin{equation}
  n(r,p) = f_\mathrm{e}-f_\mathrm{s}+1 = \left\lceil \frac{b_r}{m(\ell(p))\,\Delta} \right\rceil
\end{equation}
contiguous slots are used on every link of $p$, with $\ell(p)=\sum_{e\in p}\ell(e)$, $\Delta$ the capacity of one slot at the most robust format, and the modulation factor
\begin{equation}
  m(\delta)=
  \begin{cases}
    4 & \delta\le 500\ \text{km},\\
    3 & \delta\le 1000\ \text{km},\\
    2 & \delta\le 2000\ \text{km},\\
    1 & \text{otherwise.}
  \end{cases}
\end{equation}
\end{definition}

A request is accepted at time $t$ if such a pair exists with every slot free on every link of $p$ and no link of $p$ failed; otherwise it is blocked.
There is no queueing and no preemption: an allocation is released only at the end of its holding time or when a link of its path fails.

\subsection{Providers as overlays}

A set $\mathcal{I}=\{1,\dots,I\}$ of providers shares $G$.
Provider $i$ operates an overlay $G_i=(V_i,E_i)$ with $V_i\subseteq V$ and $E_i\subseteq E$, and every request of $i$ must be routed inside $G_i$.
Overlays overlap but decisions do not: provider $i$ never allocates spectrum on behalf of provider $j$.
The spectrum of a link is a single resource shared by all providers that use it.

\begin{definition}[Link sharing]
The number of providers that use link $e$ is
\begin{equation}
  \sigma(e) = \bigl|\{\,i\in\mathcal{I} : e\in E_i\,\}\bigr| .
\end{equation}
\end{definition}

\subsection{Disaster and timeline}

A non-empty set of nodes $V_d\subseteq\bigcap_{i}V_i$ fails at time $t_d$, together with every incident link.
Choosing $V_d$ inside the intersection of the overlays makes the threat relevant to every provider.
Provider $i$ discovers the threat at its own time $\tau_i<t_d$, without synchronization with the others, and its warning window is $[\tau_i,t_d)$.
Before $\tau_i$, it operates normally; from $\tau_i$ on, it switches to disaster-aware routing and starts migrating; after $t_d$, whatever was not migrated is lost.

\subsection{Migration}

Provider $i$ holds critical data at $s_i\in V_d$ and migrates it to a destination $d_i\in V_i\setminus V_d$.
In our instances, $d_i$ is the node of $G_i$ farthest from $V_d$ in hops, which keeps the data away from a possible extension of the damaged area.
Let $D_i$ be the volume to migrate and $M_i(t)$ the volume carried by migration requests accepted until $t$.
Migration requests originate at $s_i\in V_d$, so they must use the part of $G_i$ that is about to fail, while background traffic that must survive the failure should avoid it.
Both compete for the spectrum of the links of $G_i$ that leave $V_d$, which are also typically shared with other providers.

\subsection{Independence and cooperation}

\begin{definition}[Independent action]
Each provider is a separate agent.
There is no central allocator and no joint allocation over $\bigcup_i G_i$.
The decisions of $i$ depend only on $G_i$, on the occupancy of the links it uses and on the information it holds.
\end{definition}

\begin{definition}[Cooperation]
A subset $K\subseteq\mathcal{I}$ is \emph{cooperating} if its members share information about their overlays and migrations.
Members of $\mathcal{I}\setminus K$ keep a strictly local view.
Cooperation does not merge overlays and does not transfer the right to allocate spectrum: each $i\in K$ keeps acting only on $G_i$.
A provider joins the cooperation at its discovery time $\tau_i$.
\end{definition}

\subsection{Metrics}

Let $\mathcal{R}_\mathrm{n}$ and $\mathcal{R}_\mathrm{m}$ be the background and migration requests.
Blocking rates and availability are
\begin{equation}
  B_\mathrm{n}=\frac{|\{r\in\mathcal{R}_\mathrm{n}:\ r \text{ blocked}\}|}{|\mathcal{R}_\mathrm{n}|},\qquad
  B_\mathrm{m}=\frac{|\{r\in\mathcal{R}_\mathrm{m}:\ r \text{ blocked}\}|}{|\mathcal{R}_\mathrm{m}|},\qquad
  A=1-B .
\end{equation}
Availability here means one minus the blocking rate, not classical two-terminal availability.
The migrated fraction of provider $i$ and its mean over a set of providers $S$ are
\begin{equation}
  \eta_i=\min\!\left(1,\frac{M_i(t_d)}{D_i}\right),\qquad
  \eta(S)=\frac{1}{|S|}\sum_{i\in S}\eta_i .
\end{equation}
We report $\eta(\mathcal{I})$ for the whole network, $\eta(K)$ for the cooperating providers and $\eta(\mathcal{I}\setminus K)$ for the others, together with the background blocking of each group.

\subsection{Problem statement}

\begin{problem}[Migration in the warning window]
\label{prob:main}
\hfill
\begin{description}[leftmargin=1.6cm,style=nextline]
  \item[Given]
  a physical graph $G=(V,E)$ with $F$ slots per link;
  providers $\mathcal{I}$ with overlays $G_i$;
  a disaster set $V_d\subseteq\bigcap_i V_i$ failing at $t_d$;
  for each provider, a volume $D_i$ at $s_i\in V_d$, a destination $d_i$ and a discovery time $\tau_i<t_d$;
  background traffic on every overlay;
  contiguous and continuous spectrum allocation with distance-dependent modulation.
  \item[Find]
  for each provider, independently and from $\tau_i$ on, path and spectrum allocations in $G_i$ for its migration and its background traffic; and a cooperating subset $K$ that exchanges information.
  \item[Such that]
  the migrated fraction $\eta(\mathcal{I})$ is maximized, cooperating providers migrate at least as much as they would alone, and the background blocking $B_\mathrm{n}$ is contained.
  \item[Under]
  independent action, independent discovery, local views by default, information-only cooperation and a single spectrum per link shared by all of its users.
\end{description}
\end{problem}

%======================================================================
\section{Cooperative corridor protection}
\label{sec:mechanism}
%======================================================================

Migration and background traffic compete on the links that leave the disaster site.
A provider can favor its own migration by keeping its own background traffic out of part of the spectrum on its migration path.
If another provider also uses those links, however, the spectrum freed this way may simply be taken by the other provider's background traffic, which no one can restrict.
Protection is therefore only worthwhile on links where every user is willing to restrict itself.

\subsection{Knowledge model and exclusive links}

We assume that the infrastructure owner discloses $\sigma(e)$ for every link: a count, not the identities of the tenants.
Cooperating providers share their overlays, so for each link they know how many of \emph{them} use it.

\begin{definition}[Link exclusive to a group]
A link $e$ is exclusive to a group $S\subseteq\mathcal{I}$ if
\begin{equation}
  \bigl|\{\,i\in S : e\in E_i\,\}\bigr| = \sigma(e) > 0 ,
\end{equation}
that is, every provider that uses $e$ belongs to $S$.
\end{definition}

A single provider can prove exclusivity only for the links it uses alone ($\sigma(e)=1$).
A cooperating group can additionally prove it for links shared only among its members, and the number of such links grows with the group.

\subsection{Protected links}

The migration corridor of a provider is the set of links of its first candidate migration paths from $s_i$ to $d_i$ (its main path, or its first three paths, depending on the variant).
Whenever a provider joins the cooperation, the set of protected links is recomputed: a link is protected if it lies on the corridor of some current member of $K$ and is exclusive to $K$.
Protection is active only until $t_d$ and never applies to migration requests.
Because a protected link is used only by members of $K$, every background request that crosses it belongs to a cooperator: each cooperator restricts only its own traffic, consistently with independent action.

\subsection{Reservation size}

On a protected link, background requests may not use the top $R_e(t)$ slots, which remain available to migration.
Background requests whose only free blocks lie inside the reserved part are blocked or take another candidate path.
Two rules are studied.

\paragraph{Fixed reservation.}
$R_e = \lceil \rho F\rceil$ for a fixed fraction $\rho$, on the links of the first three migration paths.

\paragraph{Demand-driven reservation.}
Each owner $o$ of a protected link reserves, on its main path only, the slots it still needs to finish before the failure:
\begin{equation}
  R_o(t) = \left\lceil \frac{h\,\bigl(D_o-M_o(t)\bigr)}{(t_d-t)\; m_o\,\Delta} \right\rceil ,
  \qquad
  R_e(t)=\min\Bigl(F,\ \sum_{o \in \mathrm{owners}(e)} R_o(t)\Bigr),
\end{equation}
where $h$ is the mean holding time of a request and $m_o$ the modulation factor of the main path.
The reservation shrinks to zero when the owner finishes and grows if it falls behind.

\subsection{Individual protection}

To separate the effect of cooperation from the effect of knowing $\sigma(e)$, we also evaluate \emph{individual protection}: every provider that has reached its discovery time protects the links of its own corridor that it uses alone, with the same reservation rule, whether or not it cooperates.

%======================================================================
\section{Simulation setup}
\label{sec:setup}
%======================================================================

\subsection{Simulator}

We use a discrete-event simulator written in Python on SimPy and NetworkX.
Traffic is generated before the simulation and stored with the scenario, so every variant of an experiment processes exactly the same requests at the same times.
Given a scenario and a variant, the simulation is deterministic.

\subsection{Network and spectrum}

The physical topology is a 24-node, 43-link US backbone with link lengths between 250 and 2600~km (median 1000~km).
Each link has $F=200$ slots and $\Delta=12.5$~Gb/s per slot at modulation factor 1.
The modulation factor of a path follows the distance thresholds of Section~\ref{sec:problem}; with these link lengths, most paths use factor 1 or 2.
Every provider precomputes, for each pair of nodes of its overlay, the $k$ shortest paths by length ($k=10$ for disaster-aware routing).

\subsection{Providers}

$I=5$ providers are generated as random connected overlays of the topology; all overlays share at least one node and together cover every node.
In the reference scenario most links are used by one or two providers: 12 links have $\sigma(e)=1$, 21 have $\sigma(e)=2$, 4 have $\sigma(e)=3$, 5 have $\sigma(e)=4$ and one link is unused.

\subsection{Traffic}

Background requests arrive as a Poisson process over a 1200~s horizon at a rate chosen for a target network utilization of 75\%, about 156 requests per second (about $1.9\times 10^5$ requests per scenario).
Each request belongs to one provider and has its source and destination inside that provider's overlay; providers with larger and less shared overlays receive a larger share of the requests.
Bit rates are drawn uniformly from $\{100,150,\dots,400\}$~Gb/s and holding times from a normal distribution with mean 1~s and standard deviation 0.1~s.

\subsection{Disaster and migration}

The disaster is a single node chosen in the intersection of the overlays, which fails at $t_d\approx 600$~s ($\pm5\%$) and stays down for about 200~s.
Each provider discovers the threat at $\tau_i\approx 300$~s ($\pm15\%$).
Its datacenter is at the failing node and its destination is the farthest node of its overlay in hops.
From $\tau_i$ until $t_d$, the provider issues migration requests from $s_i$ to $d_i$ at a constant mean rate (about 15--17 requests per second per provider in the reference scenario), with the same bit-rate and holding-time distributions as background traffic, and the volume to migrate is $D_i\approx 4.4\times 10^5$~Gb.
Migration requests keep being issued until $t_d$ even after $D_i$ is reached, so $\eta_i$ is capped at~1.

\subsection{Routing}

Before $\tau_i$, provider $i$ routes with First-Fit over the precomputed paths of its overlay.
From $\tau_i$ on, it switches to disaster-aware First-Fit: background requests use paths that avoid $V_d$, and migration requests, which originate in $V_d$, use the shortest paths of the full overlay.
Candidate paths are tried in order of length, and the first one with a free contiguous block large enough is used.
Corridor protection, when enabled, removes the reserved slots of protected links from the blocks available to background requests.

\subsection{Compared variants}

\begin{center}
\small
\begin{tabular}{@{}lll@{}}
\toprule
Variant & Protection & Depends on $K$ \\
\midrule
First-Fit & none, not disaster-aware & no \\
No reservation (reference) & none, disaster-aware First-Fit & no \\
Fixed $\rho$, cooperative & links exclusive to $K$, $\rho\in\{10,25,40,55,70,85\}\%$ & yes \\
Fixed $\rho$, individual & links used by one provider only & no \\
Demand-driven, cooperative & links exclusive to $K$, main path & yes \\
Demand-driven, individual & links used by one provider only, main path & no \\
\bottomrule
\end{tabular}
\end{center}

\subsection{Experimental design}

For each base scenario, every variant is evaluated for all $2^I=32$ cooperating subsets $K\subseteq\mathcal{I}$, from no cooperation to full cooperation.
The order in which members join is fixed by their discovery times, so subsets and not permutations are enumerated.
Variants that do not depend on $K$ produce exactly the same simulation for every subset; they are run once and their per-provider results are relabeled for each $K$, which we verified by running one of them for both $K=\emptyset$ and $K=\mathcal{I}$.
Every variant is compared with the reference for the same scenario and the same $K$, so that differences reflect the protection mechanism only.
The results reported here use \todo{1 / 10} base scenario(s).

%======================================================================
\section{Results}
\label{sec:results}
%======================================================================

\todo{Pending the final runs without cooperative weights (Experiment\_final.ipynb). Planned content:}
\begin{itemize}[leftmargin=1.2em]
  \item \todo{Disaster-aware First-Fit vs First-Fit: migrated fraction and blocking.}
  \item \todo{Migrated fraction of all providers vs $|K|$ for each reservation rule, cooperative and individual, as gains over the reference.}
  \item \todo{Members of $K$ vs providers outside $K$: migrated fraction and change in background blocking.}
  \item \todo{Where cooperative protection overtakes individual protection, and its relation to the number of links each can prove exclusive.}
  \item \todo{Fixed vs demand-driven reservation: migration gain per point of background blocking.}
\end{itemize}

%======================================================================
\section{Limitations}
\label{sec:limitations}
%======================================================================

\begin{itemize}[leftmargin=1.2em]
  \item The tenant count $\sigma(e)$ is assumed to be disclosed by the infrastructure owner; results with $\sigma(e)$ are an upper bound for providers that would have to infer it, for example from spectrum occupancy.
  \item The disaster affects a single node; with a regional disaster, more links leaving the damaged area would be shared and fewer would be provably exclusive.
  \item Destinations are the farthest nodes of each overlay, which forces long paths with the least efficient modulation; closer destinations would need fewer slots but less distance from the threat.
  \item One topology and a single traffic model are used. \todo{number of base scenarios.}
\end{itemize}

%======================================================================
\section{Conclusion}
\label{sec:conclusion}
%======================================================================

\todo{after the results.}

\bibliographystyle{plain}
\bibliography{references}

\end{document}
